\subsection{Solving Equations with One Exponential} 
We now have a second method of solving exponential equations: by rewriting in log form. (We can think of this as taking the log of both sides!)
\begin{Example}
Solve $4^{x+1}=8$
\begin{lpic}{}{7}
\twocol{-1}{7}{Method 1:}{Method 2:};
\end{lpic}
\end{Example}
We can still use the new method if we \emph{can't} rewrite each side using the same base.
\begin{Example}
Solve $3^{x-5}=29$
\begin{lpic}{}{3}
\end{lpic}
There are two standard ways to give this answer.
\begin{itemize}
\item Give the exact solution in terms of the common log: \makeblank[1.5in]
\item Give a decimal approximation: \makeblank[1in]
\end{itemize}
\end{Example}
We may also need to \makeblank[1in] the exponential first.
\begin{Example}
Solve $\frac{1}{2}\cdot 3^{2x-1}+6=24$
\begin{lpic}{}{9}
\end{lpic}
\end{Example}